\paragraph{\trCategory。}\label{sub:categories}
取定\trCoefficients{}域 $\KK$。我们的一切\trCategory{}都默认在 $\KK$ 上\trLinear、且为小%
\trCategory。一个\textbf{\trNonunital{}\trACategory}（non-unital $A_\infty$-category）
$\scr{A}$ 由如下数据组成：一个\trObject{}的集合 $\Ob\scr{A}$；对每一对\trObject{}的一个%
\trGraded\trVectorSpace{} $\hom_{\scr{A}}(X_0,X_1)$；以及各阶 $d\ge 1$ 的\trComposition{}映射
\begin{equation}\label{eq:composition-maps}
\mu^d_{\scr{A}}\colon\hom_{\scr{A}}(X_{d-1},X_d)\otimes\cdots\otimes\hom_{\scr{A}}(X_0,X_1)
\longrightarrow\hom_{\scr{A}}(X_0,X_d)[2-d].
\end{equation}
照例，$[k]$ 表示把一个\trVectorSpace{}的\trGraded{}向下移动 $k\in\ZZ$ 位。\eqref{eq:composition-maps}
的这些映射必须满足（二次的）\textbf{\trAinfty-结合性方程}（$A_\infty$-associativity equations）
\begin{equation}\label{eq:aoo-associativity}
\sum_{m,n}(-1)^{\clubsuit_n}\,\mu_{\scr{A}}^{d-m+1}
(a_d,\ldots,a_{n+m+1},\mu_{\scr{A}}^m(a_{n+m},\ldots,a_{n+1}),a_n,\ldots,a_1)=0,
\end{equation}
其中求和取遍所有可能的项（$1\le m\le d$，$0\le n\le d-m$），而正负号由
$\clubsuit_n=|a_1|+\cdots+|a_n|-n$ 确定。同一个和还会出现在许多别的正负号表达式里，
因此我们约定，它从现在起是固定记号。

对任意\trNonunital{}\trACategory{} $\scr{A}$，都可以配上它的
\textbf{\trCohomology\trCategory}（cohomological category）$H(\scr{A})$：它与 $\scr{A}$
有相同的\trObject；\trMorphism{}空间是\trCohomology{}群
$H(\hom_{\scr{A}}(X_0,X_1),\mu^1_{\scr{A}})$；而\trComposition{}由下式定义：
\begin{equation}\label{eq:cohomological-composition}
[a_2]\cdot[a_1]=(-1)^{|a_1|}[\mu^2_{\scr{A}}(a_2,a_1)].
\end{equation}
除了可能缺少\trIdentityMorphism{}之外，这就是一个普通的\trLinear\trGraded\trCategory。
它还有一个变体，即只保留零次\trCohomology{}的子\trCategory{}%
$H^0(\scr{A})\subset H(\scr{A})$。循着同样的思路，若\trNonunital{}\trACategory{} $\scr{A}$
满足 $\mu^d_{\scr{A}}=0$（对一切 $d>2$），那么它恰就是一个
\textbf{\trNonunital{}\trDifferential\trGraded}（differential graded，dg）\trCategory。
更确切地说，在相配的 dg \trCategory{}中，$\partial a=(-1)^{|a|}\mu^1_{\scr{A}}(a)$，而
$a_2a_1=(-1)^{|a_1|}\mu^2_{\scr{A}}(a_2,a_1)$，与\eqref{eq:cohomological-composition} 一致。

\trNonunital{}\trACategory{}的\textbf{反\trCategory}（opposite category）定义如下：取
$\Ob\scr{A}^{\mathrm{opp}}=\Ob\scr{A}$ 与
$\hom_{\scr{A}^{\mathrm{opp}}}(X_0,X_1)=\hom_{\scr{A}}(X_1,X_0)$，并把\trComposition{}映射
中的顺序反过来：
\begin{equation*}
\mu^d_{\scr{A}^{\mathrm{opp}}}(a_d,\ldots,a_1)=(-1)^{\clubsuit_d}\,\mu^d_{\scr{A}}(a_1,\ldots,a_d).
\end{equation*}
在\trCohomology{}层面上，这给出通常的反\trCategory{}（不过取时不带 Koszul 正负号）：
$H(\scr{A}^{\mathrm{opp}})=H(\scr{A})^{\mathrm{opp}}$。

\begin{remark}\label{thm:terminology-conventions}
关于术语与约定的几条脚注：(i) “\trNonunital”其实应当叫“不必\trUnital{}的”。
(ii) 我们的\trCategory{}一律设定为小\trCategory，这偶尔会在形式上造成麻烦，比如定义%
\trTwistedComplex{}的时候；必要的集合论防范措施留给读者自行采取。
(iii) 使用没有\trIdentityMorphism{}的\trCategory{}时，需要小心行事：比如说，
“\trIsomorphic{}的\trObject{}”这一说法就没有合理的意义；\trFullyFaithful\trFunctor{}%
说得通，\trCategory{}之间的\trIsomorphism{}也是如此，但\trEquivalence{}则不然——至少没有
直截了当的处理方式。(iv) 我们的正负号遵循“bar 约定”，这是代数学家最常用的约定；
\eqref{eq:composition-maps} 中递减的编号，是试图把\trMorphism\trComposition{}的标准约定，与%
\trCircle{}上\trMarkedPoint{}的同样标准的循环排序调和起来，见 (9c)。
\end{remark}
